\section{Aisle management}
\label{sec:aisles}

\begin{lede}
Narrow corridors get a traffic light. Robots ask for a direction, the corridor
decides from the aggregate, and it holds that decision long enough for the
decision to be worth making --- but not forever.
\end{lede}

The design principle, stated verbatim in the code twice: \textbf{robots
generate directional requests, but aisles make directional decisions.} No
single robot can force an aisle to flip; direction is a property the aisle
decides from the aggregate of what is asking.

\subsection{Requested direction}

\code{requested\_direction(robot, aisle)} projects the robot's planned route
onto the aisle's own index ordering. If the route touches the aisle at two or
more points, it is \textsc{forward} when the index increases and
\textsc{reverse} when it decreases; for a single touching vertex it falls back
to \code{entry\_direction} (\Cref{sec:graph-aisles}).

\subsection{Per-robot directional demand}
\label{sec:aisle-demand}

\begin{equation}
  \mathrm{demand}_i(a) = w_u U_i + w_w W_i + w_p P_i - w_l L_i - w_c C_i,
  \label{eq:demand}
\end{equation}

\begin{center}\small
\begin{tabular}{@{}lll@{}}
\toprule
\textbf{Symbol} & \textbf{Definition} & \textbf{Meaning} \\
\midrule
$U_i$ & task urgency, 0 without a task & deadline pressure \\
$W_i$ & $\max(\text{waiting}, \text{no\_progress}) + \text{blocked}$ & stalled progress counts as waiting \\
$P_i$ & $1 / (1 + \min(d(x_i, \text{start}), d(x_i, \text{end})))$ & proximity to the nearer end \\
$L_i$ & route distance to waypoint, 0 if $\INF$ & long-route robots discounted \\
$C_i$ & $\mathrm{aisle\_load}(a)$ & existing congestion discourages more demand \\
\bottomrule
\end{tabular}
\end{center}

Weights: $w_u = 1.0$, $w_w = 0.5$, $w_p = 2.0$, $w_l = 0.05$, $w_c = 0.5$. The
$W_i$ term takes a maximum over two stall counters deliberately --- a robot
shuffling in place without gaining route distance is starving just as much as
one standing still, and should be able to flip the aisle.

\subsection{Aggregate demand and the decision rule}
\label{sec:aisle-decision}

\begin{align}
  S_a^{+} &= \sum_{i \ \mathrm{requesting}\ \textsc{forward}} \rho_i(a)\, \mathrm{demand}_i(a),
  &
  S_a^{-} &= \sum_{i \ \mathrm{requesting}\ \textsc{reverse}} \rho_i(a)\, \mathrm{demand}_i(a),
  \\
  \mathrm{imbalance}(a) &= S_a^{+} - S_a^{-}. &&
\end{align}

$\rho_i(a)$ is the attribution weight. By default each robot is charged to
exactly one aisle ($\rho = 1$ for its next aisle, else its current one, and $0$
elsewhere). With \param{demand\_spread} enabled the robot is charged to
\emph{every} aisle its route touches, with
$\rho_i(a) = 1 / (1 + \text{steps to the entry})$ --- the literal reading of
the specification, and the only one that lets an aisle see traffic more than
one aisle away. Measured, it engages direction control considerably more often
and costs throughput, so it ships off; the flag exists so the measurement stays
reproducible.

\code{update\_aisle\_direction} then applies hysteresis with a dead band
$[-\tau_{\mathrm{switch}}, \tau_{\mathrm{switch}}]$ ($\tau_{\mathrm{switch}} =
5.0$ when hysteresis is on, else $0$) \emph{and} a maximum green:

\begin{algorithm}[H]
\caption{\code{update\_aisle\_direction(a, t)}}
\begin{algorithmic}[1]
\If{$a$ is occupied and holds \textsc{forward}}
  \State \textbf{drain} if $\mathrm{imbalance} < -\tau_{\mathrm{switch}}$
    \textbf{ or } ($t - \mathrm{direction\_since} \geq T_{\max}$ \textbf{ and } $S_a^{-} > 0$)
\ElsIf{$a$ is occupied and holds \textsc{reverse}}
  \State \textbf{drain} if $\mathrm{imbalance} > +\tau_{\mathrm{switch}}$
    \textbf{ or } ($t - \mathrm{direction\_since} \geq T_{\max}$ \textbf{ and } $S_a^{+} > 0$)
\ElsIf{$a$ is empty and $t < \mathrm{lock\_until}$}
  \State hold the current direction \Comment{minimum-lock enforcement}
\ElsIf{$a$ is empty and the lock has expired}
  \State \textbf{commit} \textsc{forward} if $\mathrm{imbalance} > \tau_{\mathrm{switch}}$ and $S_a^{+} > 0$
  \State \textbf{commit} \textsc{reverse} if $\mathrm{imbalance} < -\tau_{\mathrm{switch}}$ and $S_a^{-} > 0$
  \State otherwise \textbf{open} the aisle
\EndIf
\end{algorithmic}
\end{algorithm}

$T_{\max} = \param{maximum\_aisle\_lock\_time}$ (default 40), floored at the
aisle's own minimum lock so the two clocks cannot contradict each other.

\paragraph{Why a maximum green is necessary.} Hysteresis is only half a traffic
signal: it bounds how \emph{soon} a committed direction may change and says
nothing about how long it may persist. Consider a warehouse with pickups down
one side and deliveries down the other --- the standard layout, and three of
the four bundled maps. Loaded robots want one direction and empty ones the
other, in near-equal measure, so $|\mathrm{imbalance}|$ sits inside the dead
band indefinitely and an aisle holds its first committed direction for the rest
of the run while half its traffic never gets through. Instrumented on
\variant{warehouse\_corridors} under a hard directional rule: all 35 robots
stalled by $t = 120$, four aisles locked \textsc{reverse} with
\code{lock\_until} long expired, the identical state still present at
$t = 199$. \Cref{prop:maxgreen} states the property the fix restores.

Flips forced this way are counted separately as \code{starvation\_flips}, which
keeps ``does hysteresis stop oscillation?'' distinct from ``does anybody
starve?''

\code{\_commit} increments \code{direction\_switches} only on an \emph{actual}
change from a non-\textsc{none} previous direction, and sets
$\mathrm{lock\_until} = t + T_{\min}$ and $\mathrm{direction\_since} = t$, so a
freshly committed direction can neither be reversed for $T_{\min}$ steps nor
held past $T_{\max}$ against opposing demand.

% worked-example: aisle
\begin{workedexample}
\textbf{Worked example.} Same run, timestep 18. Aisle 38 runs (11,20) -> (14,20) (length 4, capacity 4, axis \code{col}). Each robot routed through it contributes \code{w\_u·U + w\_w·W + w\_p·P − w\_l·L − w\_c·C} (\Cref{sec:aisle-demand}) to whichever direction it wants, and the two sides are summed:

{\fontsize{9.0}{10.8}\selectfont
\begin{verbatim}
S_a^+  (forward demand)      = -0.558
S_a^-  (reverse demand)      = 4.758
imbalance  S_a^+ − S_a^-     = −5.317
dead band  ± τ_switch        = 5
occupancy                    = 1 robot(s) inside
state on entry               = REVERSE, direction REVERSE
locked until t               = 26 (T_min = 8, T_max = 40)
\end{verbatim}
}

With hysteresis on, \code{update\_aisle\_direction} returns \code{REVERSE}; with the dead band removed it returns \code{REVERSE}. The two agree here, which is the common case: hysteresis is not meant to fire often, it is meant to make the rare flip deliberate.

An imbalance of 5.317 against a dead band of 5 clears the band, and a committed direction is additionally locked for \code{T\_min} = 8 steps. The maximum green matters for the opposite failure: past \code{T\_max} = 40 steps \emph{any} opposing demand forces a drain, so a balanced aisle - pickups on one side, deliveries on the other, imbalance permanently inside the band - cannot hold one direction forever and starve the traffic wanting the other way.
\end{workedexample}
% /worked-example

\subsection{Aisle state machine}
\label{sec:aisle-state-machine}

\begin{center}\small\ttfamily
\begin{tabular}{@{}rcl@{}}
\toprule
\multicolumn{3}{@{}l@{}}{\textsf{imbalance past $\pm\tau$, \textbf{or} held $\geq T_{\max}$ with opposing demand}}\\
\textsc{forward} & $\longrightarrow$ & \textsc{draining} \\
\textsc{reverse} & $\longrightarrow$ & \textsc{draining} \\
\midrule
\textsc{draining} & $\longrightarrow$ & \textsc{reverse}/\textsc{forward} \ \textsf{(occupancy reaches 0: commit pending)}\\
\textsc{draining} & $\longrightarrow$ & \textsc{open} \ \textsf{(drain timed out)}\\
\textsc{open} & $\longrightarrow$ & \textsc{forward}/\textsc{reverse} \ \textsf{(imbalance + lock expired)}\\
\bottomrule
\end{tabular}
\end{center}

Entering \textsc{draining} records a \code{pending\_direction}: the opposite of
the current one. When the aisle empties, that direction is committed
\textbf{directly}, without re-running the imbalance test. The demand that
triggered the drain has usually moved on by the time the aisle is clear, so
re-testing would simply re-commit the direction just drained, and the flip
would never happen.

One addition beyond the base rule: a \textsc{draining} aisle that has not
reached zero occupancy within \param{max\_drain\_time} (default 30) steps is
force-reopened. The specification allows an immediate switch when ``the current
direction is infeasible'', and an aisle that cannot drain is exactly that.
Without it the state machine has an absorbing deadlock.

\subsubsection{Coordinated direction assignment}
\label{sec:aisle-coordination}

\param{coordinate\_aisle\_directions} (default off) decides directions as a
\emph{set} rather than one aisle at a time: aisles commit in descending
$|\mathrm{imbalance}|$, and any commit that would break strong connectivity of
the directed residual graph is rolled back, leaving that aisle bidirectional.
The check is a two-way BFS over passable cells under the current assignment,
$O(|V| + |E|)$ per commit.

\begin{caveat}
This was the fix an earlier review expected to matter most. It does not: on the
bundled maps its measured effect on throughput is between $\pm 0.000$ and
$\pm 0.011$, because a single one-way aisle almost never disconnects a ladder
graph, so the guard hardly ever fires. The collapse it was meant to cure was a
liveness failure (\Cref{sec:aisle-decision}), not a connectivity failure. It
ships as a flag so the null result stays reproducible.
\end{caveat}

\subsection{Reservations}
\label{sec:aisle-reservations}

\code{update\_aisle\_reservations} expires any reservation past its
\code{expiry\_time}, then grants a fresh one only if \emph{all} of: the aisle
is not \textsc{draining}; the robot's requested direction matches the aisle's
current one, or the aisle has none; the robot does not already hold one;
occupancy plus pending reservations held by robots still outside is below
capacity; and the robot is within route distance 3 of either endpoint. A
granted reservation carries
\begin{equation}
  \mathrm{expected\_exit} = t + \lfloor \ell \rfloor + 1,
  \qquad
  \mathrm{expiry} = t + \param{reservation\_ttl} \ (= 15).
\end{equation}

The read side, \code{can\_enter\_aisle}, splits into two gates because they
answer different questions:

\begin{itemize}
\item the \textbf{occupancy cap} applies to every managed aisle,
  \textsc{open} included. Over-filling a single-file corridor is what builds a
  queue that cannot drain, and that is true whether or not the corridor
  currently has a direction. This is the mechanism hypothesis H3 is about.
\item the \textbf{ticket} is required only once the aisle is directional.
  Demanding one to enter an \textsc{open} aisle throttles every crossing on a
  map whose aisles are mostly short and mostly open.
\end{itemize}

\begin{equation}
  \mathrm{can\_enter}(i, a) =
  \begin{cases}
    \text{false} & a \text{ is \textsc{draining}},\\
    \text{false} & a \text{ is directional, direction mismatched},\\
    \text{false} & \mathrm{occupancy}(a) \geq \mathrm{capacity}(a),\\
    \text{true} & \text{reservations off, or } a \text{ is \textsc{open}},\\
    \mathrm{has\_ticket}(i, a) & \text{otherwise.}
  \end{cases}
\end{equation}

\begin{caveat}
\textbf{Gating.} These are gated on \param{reservations} alone. They used to be
gated on \code{AisleManager.enabled} (that is, on
\param{direction\_control} $=$ \code{"aisle"}), which made the flag a silent
no-op in every robot-level configuration --- including \variant{reservations\_only},
the variant that exists specifically to isolate it, whose runs were
bit-identical to its own control. An entire factorial design was built on that
empty cell (\Cref{sec:results-factorial}).
\end{caveat}

\subsection{Movement legality}
\label{sec:aisle-legality}

One deliberate deviation: direction restricts \textbf{entry}, not
\textbf{interior movement}. A robot already inside a directional aisle may
always make an \emph{egress} move toward the nearer endpoint. With
$c = \mathrm{index}(\text{current})$, $n = \mathrm{index}(\text{candidate})$
and $\ell$ the last index,
\begin{equation}
  \mathrm{is\_egress}(c, n) \iff
    \min(n, \ell - n) < \min(c, \ell - c),
\end{equation}
that is, the candidate is strictly closer in index-distance to \emph{whichever}
endpoint is nearest, regardless of which one that is. Without this, a
\textsc{draining} aisle could trap a robot facing the ``wrong'' way, and the
aisle would never empty.
